%!TEX root = ../main.tex
\section{Proof of Theorem \ref{thm:nonlinear-bound} }
\label{sec:app-proof-nonlinear-upperbound}

\begin{proof}
We prove \eqref{eq:E2_theta} and \eqref{eq:E2_ell} separately.

\paragraph{Part (i): proof of \eqref{eq:E2_theta}.}
Let
\[
u(\tau):=\sum_{t=0}^{T-1} J_\theta(s_t)^\top \Sigma^{-1}\varepsilon_t,
\qquad\text{so that}\qquad
\widehat G_\theta(\tau)=R(\tau)\,u(\tau).
\]
\emph{Step 1 (sum-of-norms).}
Using $\big\|\sum_{t=0}^{T-1} x_t\big\|_2^2 \le T\sum_{t=0}^{T-1}\|x_t\|_2^2$,
with $x_t=J_\theta(s_t)^\top \Sigma^{-1}\varepsilon_t$, we have
\[
\|u(\tau)\|_2^2 \le T\sum_{t=0}^{T-1}\|J_\theta(s_t)^\top \Sigma^{-1}\varepsilon_t\|_2^2.
\]
Hence
\begin{align}
\mathbb E\!\left[\|\widehat G_\theta(\tau)\|_2^2 \,\middle|\, s_0\right]
&=
\mathbb E\!\left[R(\tau)^2\|u(\tau)\|_2^2 \,\middle|\, s_0\right] \nonumber\\
&\le
T\sum_{t=0}^{T-1}
\mathbb E\!\left[R(\tau)^2\|J_\theta(s_t)^\top \Sigma^{-1}\varepsilon_t\|_2^2 \,\middle|\, s_0\right].
\label{eq:step1}
\end{align}

\emph{Step 2 (operator/Frobenius bound).}
For each $t$,
\[
\|J_\theta(s_t)^\top \Sigma^{-1}\varepsilon_t\|_2
\le \|J_\theta(s_t)\|_{\mathrm{F}}\,\|\Sigma^{-1}\|_2\,\|\varepsilon_t\|_2,
\]
Plugging into \eqref{eq:step1} gives
\begin{equation}
\label{eq:step2}
\mathbb E\!\left[\|\widehat G_\theta(\tau)\|_2^2 \,\middle|\, s_0\right]
\le
T\|\Sigma^{-1}\|_2^2\sum_{t=0}^{T-1}
\mathbb E\!\left[R(\tau)^2\|J_\theta(s_t)\|_{\mathrm{F}}^2\|\varepsilon_t\|_2^2 \,\middle|\, s_0\right].
\end{equation}

\emph{Step 3 (Cauchy--Schwarz).}
For each $t$, apply Cauchy--Schwarz with
$X=R(\tau)^2$ and $Y=\|J_\theta(s_t)\|_{\mathrm{F}}^2\|\varepsilon_t\|_2^2$:
\[
\mathbb E[XY\mid s_0]\le \sqrt{\mathbb E[X^2\mid s_0]\,\mathbb E[Y^2\mid s_0]}
=
\sqrt{\mathbb E[R(\tau)^4\mid s_0]}\,
\sqrt{\mathbb E[\|J_\theta(s_t)\|_{\mathrm{F}}^4\|\varepsilon_t\|_2^4\mid s_0]}.
\]
Thus \eqref{eq:step2} implies
\begin{equation}
\label{eq:step3}
\mathbb E\!\left[\|\widehat G_\theta(\tau)\|_2^2 \,\middle|\, s_0\right]
\le
T\|\Sigma^{-1}\|_2^2\sqrt{\mathbb E[R(\tau)^4\mid s_0]}
\sum_{t=0}^{T-1}
\sqrt{\mathbb E[\|J_\theta(s_t)\|_{\mathrm{F}}^4\|\varepsilon_t\|_2^4\mid s_0]}.
\end{equation}

\emph{Step 4 (independence of $\varepsilon_t$ and $s_t$).}
Since $s_t$ is a measurable function of $(s_0,\varepsilon_0,\dots,\varepsilon_{t-1})$
and $\varepsilon_t$ is independent of $(\varepsilon_0,\dots,\varepsilon_{t-1})$,
we have $\varepsilon_t\perp s_t$ conditional on $s_0$. Therefore
\[
\mathbb E[\|J_\theta(s_t)\|_{\mathrm{F}}^4\|\varepsilon_t\|_2^4\mid s_0]
=
\mathbb E[\|J_\theta(s_t)\|_{\mathrm{F}}^4\mid s_0]\;\mathbb E[\|\varepsilon\|_2^4].
\]

\emph{Step 5 (Gaussian fourth moment).}
For $\varepsilon\sim\mathcal N(0,\Sigma)$,
\[
\mathbb E\|\varepsilon\|_2^4
=
(\operatorname{tr}\Sigma)^2 + 2\,\operatorname{tr}(\Sigma^2).
\]
Substituting the last two displays into \eqref{eq:step3} yields \eqref{eq:E2_theta}.
\paragraph{Part (ii): proof of \eqref{eq:E2_ell}.}
Define
\[
q(\varepsilon_t):=\nabla_\ell \log\pi(a_t\mid s_t)
=\Sigma^{-1}(\varepsilon_t\odot\varepsilon_t)-\mathbf 1
=\big((\varepsilon_t\odot\varepsilon_t)\odot e^{-2\ell}\big)-\mathbf 1\in\mathbb R^m,
\]
so that
\[
\widehat G_\ell(\tau)=R(\tau)\sum_{t=0}^{T-1} q(\varepsilon_t).
\]

\emph{Step 1 (sum-of-norms).}
Using $\big\|\sum_{t=0}^{T-1} x_t\big\|_2^2 \le T\sum_{t=0}^{T-1}\|x_t\|_2^2$ with $x_t=q(\varepsilon_t)$:
\[
\|\widehat G_\ell(\tau)\|_2^2
=
R(\tau)^2\Big\|\sum_{t=0}^{T-1}q(\varepsilon_t)\Big\|_2^2
\le
R(\tau)^2\,T\sum_{t=0}^{T-1}\|q(\varepsilon_t)\|_2^2.
\]
Taking $\mathbb E[\cdot\mid s_0]$ and applying Cauchy--Schwarz termwise gives
\begin{align}
\mathbb E[\|\widehat G_\ell(\tau)\|_2^2\mid s_0]
&\le
T\sum_{t=0}^{T-1}
\mathbb E\!\left[R(\tau)^2\|q(\varepsilon_t)\|_2^2\mid s_0\right]\nonumber\\
&\le
T\sum_{t=0}^{T-1}
\sqrt{\mathbb E\!\left[R(\tau)^4\mid s_0\right]}\;
\sqrt{\mathbb E\!\left[\|q(\varepsilon_t)\|_2^4\right]} \nonumber\\
&=
T^2\sqrt{\mathbb E\!\left[R(\tau)^4\mid s_0\right]}\;
\sqrt{\mathbb E\!\left[\|q(\varepsilon)\|_2^4\right]},
\label{eq:ell_bound_reduce}
\end{align}
where the last line uses that $\varepsilon_t$ are i.i.d.

\emph{Step 2 (exact fourth moment of the log-std score).}
Assume $\ell=\log\sigma$ and $\Sigma=\diag{e^{2\ell}}$ is diagonal. Then $\varepsilon=\diag{e^{\ell}}z$ with $z\sim\mathcal N(0,I_m)$, and elementwise
\[
q_i(\varepsilon)=\frac{\varepsilon_i^2}{e^{2\ell_i}}-1=z_i^2-1,
\]
so the distribution of $q(\varepsilon)$ does not depend on $\ell$.
Let $w_i:=(z_i^2-1)^2$. Then
\[
\|q(\varepsilon)\|_2^2=\sum_{i=1}^m w_i,
\qquad
\|q(\varepsilon)\|_2^4=\left(\sum_{i=1}^m w_i\right)^2
=\sum_{i=1}^m w_i^2 + 2\sum_{1\le i<j\le m} w_iw_j.
\]
Using independence of $(z_i)$, we have $\mathbb E[w_iw_j]=\mathbb E[w_i]\mathbb E[w_j]$ for $i\neq j$.
Moreover, for $z\sim\mathcal N(0,1)$,
\[
\mathbb E[(z^2-1)^2]=2,\qquad \mathbb E[(z^2-1)^4]=60.
\]
Therefore
\[
\mathbb E\|q(\varepsilon)\|_2^4
=
m\cdot 60 + 2\binom{m}{2}\cdot (2\cdot 2)
=
4m(m+14).
\]
Plugging this into \eqref{eq:ell_bound_reduce} yields
\[
\mathbb E[\|\widehat G_\ell(\tau)\|_2^2\mid s_0]
\le
2T^2\sqrt{m(m+14)}\;\sqrt{\mathbb E[R(\tau)^4\mid s_0]}.
\]

\end{proof}
